\subsection{Navegación autónoma y mapeo con LiDAR simulado}

\subsubsection{Descripción de la prueba}

Esta prueba valida el comportamiento del módulo de navegación autónoma en un entorno virtual con mapeo simultáneo. El robot debe construir un mapa de ocupación a partir del LiDAR simulado, planificar trayectoria y completar un recorrido prolongado sin perder la consistencia global de localización.

El ensayo incluye además un evento dinámico: un objetivo cruza la trayectoria del robot durante la navegación. Frente a esta condición, el sistema debe interrumpir el avance, priorizar el seguimiento del objetivo y recién retomar la marcha cuando el evento finaliza.

\paragraph{Objetivos mínimos}
\begin{itemize}
    \item Verificar que el sistema completa una trayectoria equivalente a 200 veces la longitud del robot.
    \item Validar la calidad del mapa generado respecto del mapa de referencia del entorno.
    \item Verificar el error de posicionamiento global bajo modelado de medición con ruido tipo GPS.
    \item Verificar la conmutación de estado de \texttt{NAVIGATING} a \texttt{TRACKING} ante cruce de objetivo.
    \item Medir tiempos de ejecución y continuidad de operación sin colisiones.
\end{itemize}

\subsubsection{Criterios de aprobación}

Tomando una longitud de robot de $L_r = 0{,}60\ \text{m}$, la distancia mínima de prueba queda definida como:
\[
D_{\min} = 200 \cdot L_r = 120\ \text{m}
\]

\begin{table}[H]
\centering
\small
\begin{tabular}{p{8.8cm} p{4.1cm}}
\toprule
\textbf{Métrica} & \textbf{Criterio de aprobación} \\
\midrule
Distancia total recorrida & $\geq 120\ \text{m}$ \\
Tiempo total para completar el recorrido & $\leq 10\ \text{min}$ \\
Colisiones con obstáculos estáticos o dinámicos & 0 eventos \\
Error global de posicionamiento medio (trayectoria estimada vs referencia) & $\leq 3{,}5\ \text{m}$ \\
Percentil 95 del error global de posicionamiento & $\leq 5{,}0\ \text{m}$ \\
Calidad del mapa de ocupación (IoU mapa generado vs mapa de referencia) & $\geq 0{,}75$ \\
Cobertura del área navegable en el mapa generado & $\geq 90\%$ \\
Tiempo de reacción ante cruce de objetivo (detección a detención) & $\leq 1{,}5\ \text{s}$ \\
Avance residual luego de detectar objetivo en cruce & $\leq 0{,}50\ \text{m}$ \\
\bottomrule
\end{tabular}
\caption{Criterios de aprobación de la prueba de navegación y mapeo en simulación.}
\label{tab:criterios_prueba_navegacion}
\end{table}

\subsubsection{Procedimiento de ejecución}

\begin{enumerate}
    \item Configurar escenario virtual con obstáculos estáticos, al menos un corredor estrecho y un tramo abierto para navegación prolongada.
    \item Inicializar el robot en estado \texttt{NAVIGATING}, con LiDAR simulado activo y módulo de mapeo habilitado.
    \item Definir ruta de misión para asegurar distancia acumulada objetivo de $120\ \text{m}$ o superior.
    \item Activar registro de datos: trayectoria de referencia, trayectoria estimada, mapa generado, tiempos de estado y eventos de detección.
    \item Inyectar ruido en posición global reportada para emular medición tipo GPS con precisión de $\pm 5\ \text{m}$.
    \item Durante el recorrido, forzar el cruce transversal de un objetivo en la trayectoria del robot.
    \item Verificar transición de estados hacia \texttt{TRACKING}, detención del avance y posterior retorno a \texttt{NAVIGATING}.
    \item Finalizar el ensayo al superar la distancia mínima o al alcanzar condición de falla.
\end{enumerate}

\subsubsection{Resultados}

Los siguientes campos se completan luego de ejecutar la simulación.

\begin{table}[H]
\centering
\small
\begin{tabular}{p{8.8cm} p{2.2cm} p{2.8cm}}
\toprule
\textbf{Métrica} & \textbf{Resultado} & \textbf{Cumple} \\
\midrule
Distancia total recorrida [m] & --- & --- \\
Tiempo total de recorrido [min:s] & --- & --- \\
Cantidad de colisiones [\#] & --- & --- \\
Error global medio de posicionamiento [m] & --- & --- \\
Percentil 95 del error global [m] & --- & --- \\
IoU mapa generado vs referencia [-] & --- & --- \\
Cobertura del mapa navegable [\%] & --- & --- \\
Tiempo detección-detención ante cruce [s] & --- & --- \\
Avance residual post-detección [m] & --- & --- \\
\bottomrule
\end{tabular}
\caption{Resultados medidos de la prueba de navegación y mapeo en simulación (a completar).}
\label{tab:resultados_prueba_navegacion}
\end{table}

De forma complementaria, se recomienda incluir en la versión final del informe:
\begin{itemize}
    \item Mapa de referencia y mapa generado superpuestos.
    \item Trayectoria de referencia y trayectoria estimada en el mismo plano.
    \item Registro temporal de estados de misión (\texttt{NAVIGATING}/\texttt{TRACKING}).
\end{itemize}
